\chapter{Comandos de programación}
\label{cp:scpi-commands}

Este capítulo describe el conjunto de comandos SCPI (Standard Commands for
Programmable Instruments) implementados en el \textit{firmware}. Los comandos
se agrupan en los subsistemas estándar \texttt{MEASure}, \texttt{SENSe},
\texttt{SOURce} y \texttt{OUTPut}, a los que se añaden los comandos comunes
definidos en el estándar IEEE~488.2.

\section{Convenciones de notación}

\begin{itemize}
    \item El carácter \texttt{\#} en un mnemónico representa un sufijo numérico
          que identifica el canal (1 o 2). Por ejemplo,
          \texttt{MEASure:TEMPerature\#?} acepta tanto
          \texttt{MEAS:TEMP1?} como \texttt{MEAS:TEMP2?}.
    \item Las palabras clave se muestran en mayúsculas para la parte abreviada
          y en minúsculas para el resto. La forma abreviada es la que reconoce
          el \textit{parser}.
    \item Los parámetros entre corchetes (\texttt{[ ]}) son opcionales. Si se
          omiten, el comando se ejecuta con el valor predeterminado.
    \item Los ejemplos de sesión muestran el comando enviado por el
          \textit{host} precedido por \texttt{>} y la respuesta del
          instrumento precedida por \texttt{<}. El carácter \texttt{\\n}
          representa el terminador de mensaje.
\end{itemize}

\section{Comandos comunes IEEE~488.2}

\subsection{Identificación}

\begin{description}
    \item[Sintaxis:] \texttt{*IDN?\n}
    \item[Descripción:] Retorna la cadena de identificación del instrumento en
          el formato \texttt{<fabricante>,<modelo>,<número de serie>,<versión>}.
    \item[Respuesta:] \texttt{DF,CONTROL DE TEMPERATURA,CT2CH,v1.1}
    \item[Ejemplo:]
          \begin{verbatim}
          > *IDN?\n
          < DF,CONTROL DE TEMPERATURA,CT2CH,v1.1\n
          \end{verbatim}
\end{description}

\subsection{Reset}

\begin{description}
    \item[Sintaxis:] \texttt{*RST\n}
    \item[Descripción:] Restaura el instrumento a sus valores por defecto. Apaga
          todas las salidas, carga los parámetros de fábrica del PID y el
          \textit{setpoint}, y limpia los registros de estado IEEE~488.2.
    \item[Ejemplo:]
          \begin{verbatim}
          > *RST\n
          \end{verbatim}
\end{description}

\subsection{Clear Status}

\begin{description}
    \item[Sintaxis:] \texttt{*CLS\n}
    \item[Descripción:] Borra todos los registros de evento y la cola de
          errores.
    \item[Ejemplo:]
          \begin{verbatim}
          > *CLS\n
          \end{verbatim}
\end{description}

\subsection{Autocomprobación}

\begin{description}
    \item[Sintaxis:] \texttt{*TST?\n}
    \item[Descripción:] Ejecuta una autocomprobación y retorna el resultado.
    \item[Respuesta:] \texttt{0} (sin errores) o \texttt{1} (fallo detectado).
    \item[Ejemplo:]
          \begin{verbatim}
          > *TST?\n
          < 0\n
          \end{verbatim}
\end{description}

\section{Subsistema \texttt{MEASure}}

\subsection{Medición de temperatura}

\begin{description}
    \item[Sintaxis:] \texttt{MEASure:TEMPerature\#?\n}
    \item[Parámetros:] \texttt{\#} = número de sensor (1 o 2).
    \item[Descripción:] Retorna la temperatura medida por el sensor indicado, en
          grados Celsius.
    \item[Respuesta:] \texttt{<float>}~°C
    \item[Ejemplo:]
          \begin{verbatim}
          > MEAS:TEMP1?\n
          < 24.7\n
          \end{verbatim}
\end{description}

\section{Subsistema \texttt{SENSe}}

\subsection{Tipo de sensor}

\begin{description}
    \item[Sintaxis:] \texttt{SENSe\#:TEMPerature:TYPE <tipo>\n}
          \texttt{SENSe\#:TEMPerature:TYPE?\n}
    \item[Parámetros:] \texttt{\#} = número de sensor (1 o 2).\\
          \texttt{<tipo>} = \texttt{RTD}, \texttt{TCouple}, \texttt{LM35},
          \texttt{NTC}.
    \item[Descripción:] Configura o consulta el tipo de elemento sensor
          conectado al canal.
    \item[Ejemplo:]
          \begin{verbatim}
          > SENS1:TEMP:TYPE RTD\n
          > SENS1:TEMP:TYPE?\n
          < RTD\n
          \end{verbatim}
\end{description}

\subsection{Protección del sensor}

\subsubsection{Límite inferior}

\begin{description}
    \item[Sintaxis:] \texttt{SENSe\#:TEMPerature:PROTection:MIN <valor>\n}
          \texttt{SENSe\#:TEMPerature:PROTection:MIN?\n}
    \item[Parámetros:] \texttt{<valor>} = temperatura en °C.
    \item[Descripción:] Establece o consulta el límite inferior de protección.
          Si la temperatura medida desciende por debajo de este valor, la salida
          correspondiente se desactiva.
    \item[Ejemplo:]
          \begin{verbatim}
          > SENS1:TEMP:PROT:MIN -10\n
          > SENS1:TEMP:PROT:MIN?\n
          < -10.0\n
          \end{verbatim}
\end{description}

\subsubsection{Límite superior}

\begin{description}
    \item[Sintaxis:] \texttt{SENSe\#:TEMPerature:PROTection:MAX <valor>\n}
          \texttt{SENSe\#:TEMPerature:PROTection:MAX?\n}
    \item[Parámetros:] \texttt{<valor>} = temperatura en °C.
    \item[Descripción:] Establece o consulta el límite superior de protección.
          Si la temperatura medida supera este valor, la salida correspondiente
          se desactiva.
    \item[Ejemplo:]
          \begin{verbatim}
          > SENS1:TEMP:PROT:MAX 120\n
          > SENS1:TEMP:PROT:MAX?\n
          < 120.0\n
          \end{verbatim}
\end{description}

\subsubsection{Estado de la protección}

\begin{description}
    \item[Sintaxis:] \texttt{SENSe\#:TEMPerature:PROTection[:STATe] <ON|OFF>\n}
          \texttt{SENSe\#:TEMPerature:PROTection[:STATe]?\n}
    \item[Parámetros:] \texttt{ON} (protección habilitada) u \texttt{OFF}
          (protección deshabilitada).
    \item[Descripción:] Activa, desactiva o consulta la protección de
          temperatura en el canal indicado.
    \item[Ejemplo:]
          \begin{verbatim}
          > SENS1:TEMP:PROT ON\n
          > SENS1:TEMP:PROT?\n
          < ON\n
          \end{verbatim}
\end{description}

\subsubsection{Protección disparada}

\begin{description}
    \item[Sintaxis:] \texttt{SENSe\#:TEMPerature:PROTection:TRIPped?\n}
    \item[Descripción:] Retorna \texttt{ON} si la protección del canal está
          actualmente disparada (la temperatura ha excedido un límite y no se
          ha rearmado), \texttt{OFF} en caso contrario.
    \item[Ejemplo:]
          \begin{verbatim}
          > SENS1:TEMP:PROT:TRIP?\n
          < OFF\n
          \end{verbatim}
\end{description}

\section{Subsistema \texttt{SOURce}}

\subsection{Sensor de realimentación}

\begin{description}
    \item[Sintaxis:] \texttt{SOURce\#:CONTrol:FEEDback <sensor>\n}
          \texttt{SOURce\#:CONTrol:FEEDback?\n}
    \item[Parámetros:] \texttt{\#} = número de lazo de control (1 o 2).\\
          \texttt{<sensor>} = número de sensor (1 o 2) que proporciona la
          temperatura de entrada al lazo.
    \item[Descripción:] Asigna o consulta qué sensor físico alimenta el lazo de
          control indicado.
    \item[Ejemplo:]
          \begin{verbatim}
          > SOUR1:CONT:FEED 1\n
          > SOUR1:CONT:FEED?\n
          < 1\n
          \end{verbatim}
\end{description}

\subsection{Punto de consigna}

\begin{description}
    \item[Sintaxis:] \texttt{SOURce\#:CONTrol:SETPoint <valor>\n}
          \texttt{SOURce\#:CONTrol:SETPoint?\n}
    \item[Parámetros:] \texttt{<valor>} = temperatura en °C (rango: -10 a 120).
    \item[Descripción:] Establece o consulta el punto de consigna del lazo de
          control.
    \item[Ejemplo:]
          \begin{verbatim}
          > SOUR1:CONT:SETP 35.5\n
          > SOUR1:CONT:SETP?\n
          < 35.5\n
          \end{verbatim}
\end{description}

\subsection{Parámetros PID}

\subsubsection{Ganancia proporcional}

\begin{description}
    \item[Sintaxis:] \texttt{SOURce\#:CONTrol:PID:KP <valor>\n}
          \texttt{SOURce\#:CONTrol:PID:KP?\n}
    \item[Parámetros:] \texttt{<valor>} = ganancia proporcional ($K_p$).
    \item[Descripción:] Establece o consulta la ganancia proporcional del PID.
    \item[Ejemplo:]
          \begin{verbatim}
          > SOUR1:CONT:PID:KP 2.5\n
          > SOUR1:CONT:PID:KP?\n
          < 2.5000\n
          \end{verbatim}
\end{description}

\subsubsection{Ganancia integral}

\begin{description}
    \item[Sintaxis:] \texttt{SOURce\#:CONTrol:PID:KI <valor>\n}
          \texttt{SOURce\#:CONTrol:PID:KI?\n}
    \item[Parámetros:] \texttt{<valor>} = ganancia integral ($K_i$).
    \item[Descripción:] Establece o consulta la ganancia integral del PID.
    \item[Ejemplo:]
          \begin{verbatim}
          > SOUR1:CONT:PID:KI 0.05\n
          > SOUR1:CONT:PID:KI?\n
          < 0.0500\n
          \end{verbatim}
\end{description}

\subsubsection{Ganancia derivativa}

\begin{description}
    \item[Sintaxis:] \texttt{SOURce\#:CONTrol:PID:KD <valor>\n}
          \texttt{SOURce\#:CONTrol:PID:KD?\n}
    \item[Parámetros:] \texttt{<valor>} = ganancia derivativa ($K_d$).
    \item[Descripción:] Establece o consulta la ganancia derivativa del PID.
    \item[Ejemplo:]
          \begin{verbatim}
          > SOUR1:CONT:PID:KD 0.01\n
          > SOUR1:CONT:PID:KD?\n
          < 0.0100\n
          \end{verbatim}
\end{description}

\subsection{Límites de salida}

\subsubsection{Límite inferior}

\begin{description}
    \item[Sintaxis:] \texttt{SOURce\#:CONTrol:OUTPut:LIMit:MIN <valor>\n}
          \texttt{SOURce\#:CONTrol:OUTPut:LIMit:MIN?\n}
    \item[Parámetros:] \texttt{<valor>} = ciclo de trabajo mínimo (0 a 100).
    \item[Descripción:] Establece o consulta el límite inferior de la salida del
          PID.
    \item[Ejemplo:]
          \begin{verbatim}
          > SOUR1:CONT:OUTP:LIM:MIN 0\n
          > SOUR1:CONT:OUTP:LIM:MIN?\n
          < 0.0\n
          \end{verbatim}
\end{description}

\subsubsection{Límite superior}

\begin{description}
    \item[Sintaxis:] \texttt{SOURce\#:CONTrol:OUTPut:LIMit:MAX <valor>\n}
          \texttt{SOURce\#:CONTrol:OUTPut:LIMit:MAX?\n}
    \item[Parámetros:] \texttt{<valor>} = ciclo de trabajo máximo (0 a 100).
    \item[Descripción:] Establece o consulta el límite superior de la salida del
          PID.
    \item[Ejemplo:]
          \begin{verbatim}
          > SOUR1:CONT:OUTP:LIM:MAX 80\n
          > SOUR1:CONT:OUTP:LIM:MAX?\n
          < 80.0\n
          \end{verbatim}
\end{description}

\subsubsection{Ciclo de trabajo actual}

\begin{description}
    \item[Sintaxis:] \texttt{SOURce\#:CONTrol:OUTPut:LIMit:STATe?\n}
    \item[Descripción:] Retorna el ciclo de trabajo actual de la salida del PID
          en porcentaje.
    \item[Ejemplo:]
          \begin{verbatim}
          > SOUR1:CONT:OUTP:LIM:STAT?\n
          < 45.3\n
          \end{verbatim}
\end{description}

\section{Subsistema \texttt{OUTPut}}

\subsection{Estado de la salida física}

\begin{description}
    \item[Sintaxis:] \texttt{OUTPut\#:STATe <ON|OFF>\n}
          \texttt{OUTPut\#:STATe?\n}
    \item[Parámetros:] \texttt{ON} (salida activada) u \texttt{OFF} (salida
          desactivada).
    \item[Descripción:] Activa, desactiva o consulta el estado de la etapa de
          potencia del canal indicado. La protección del sensor tiene prioridad:
          aunque se ordene \texttt{ON}, la salida no se activará si la
          protección correspondiente está disparada.
    \item[Ejemplo:]
          \begin{verbatim}
          > OUTP1:STAT ON\n
          > OUTP1:STAT?\n
          < ON\n
          \end{verbatim}
\end{description}